\documentclass[12pt,a4paper]{ctexart}
\usepackage{geometry}
\geometry{left=2.5cm,right=2.5cm,top=2.5cm,bottom=2.5cm}
\usepackage{amsmath, amssymb}
\usepackage{tabularx}
\usepackage{longtable}
\usepackage{hyperref}
\usepackage{enumitem}

\title{AI 工具使用详情}
\author{}
\date{2026-08-21 到 2026-08-22}

\begin{document}
\maketitle

% ============================================================
% 一、所用 AI 工具名称和版本
% ============================================================
\section*{一、所用 AI 工具名称和版本}

\begin{longtable}{|c|l|l|l|l|}
\hline
序号 & AI工具名称 & 版本/型号 & 开发机构/公司 & 使用日期 \\
\hline
[AI-1-1] & Trae & 1.0（CN） & 字节跳动 & 2026-08-21 到 2026-08-22 \\
\hline
[AI-1-2] & DeepSeek & DeepSeek-V4 & 深度求索（DeepSeek） & 2026-08-21 到 2026-08-22 \\
\hline
\end{longtable}

% ============================================================
% 二、具体使用目的和环节
% ============================================================
\section*{二、具体使用目的和环节}

\begin{longtable}{|c|p{3.2cm}|p{6.2cm}|p{2.6cm}|}
\hline
序号 & 使用环节（论文章节） & 使用目的 & 使用的AI工具 \\
\hline
[AI-1-1] & 赛题背景与领域知识（阶段 0） & 三疾病宏基因组文献检索、数据来源与已知标志物整理、参考解答拉取 & DeepSeek \\
\hline
[AI-1-2] & 第 N 章 疾病预测模型（S1 方案设计） & 题型判定、模型选型（L2+CLR 主模型 / RF 对照 / 基线）、评估协议确定 & DeepSeek \\
\hline
[AI-1-3] & 第 N 章 疾病预测模型（S1 建模实现） & 近全零过滤、CLR 变换、Logistic(L2)、随机森林、基线、分层 5 折 CV 与 LOOCV 代码实现 & DeepSeek, Trae \\
\hline
[AI-1-4] & 第 N 章 疾病预测模型（S1 结果分析） & 三数据集性能对比、跨疾病差异四重归因、待裁定项裁决 & DeepSeek \\
\hline
[AI-2-1] & 第 N 章 特征选择与生物标志物（S2 方案设计） & 稳定性选择（Lasso+bootstrap）、两路信号、共现分析、VIP 佐证方案确定 & DeepSeek \\
\hline
[AI-2-2] & 第 N 章 特征选择与生物标志物（S2 建模实现） & 近全零过滤、CLR、Lasso+bootstrap、Fisher/Wilcoxon 检验、共现分析、VIP 代码实现 & DeepSeek, Trae \\
\hline
[AI-2-3] & 第 N 章 特征选择与生物标志物（S2 结果分析） & 每病稳定标志物、两路信号、共现、跨疾病对比、待裁定项裁决 & DeepSeek \\
\hline
[AI-3-1] & 第 N 章 跨疾病预测模型（S3 方案设计） & LODO 协议、四策略对比、紧急回退协议确定 & DeepSeek \\
\hline
[AI-3-2] & 第 N 章 跨疾病预测模型（S3 建模实现） & 四策略、回退 R1--R4、衰减归因、阈值漂移代码实现 & DeepSeek, Trae \\
\hline
[AI-3-3] & 第 N 章 跨疾病预测模型（S3 结果分析） & 四策略对比、回退、衰减归因、深度迁移、C3 阈值漂移分析 & DeepSeek \\
\hline
\end{longtable}

% ============================================================
% 三、关键交互记录
% ============================================================
\section*{三、关键交互记录}

\begin{description}[leftmargin=*, itemindent=0pt, labelwidth=*, labelsep=0.5em]

\item[\textbf{交互记录 [AI-1-1] —— 领域知识整理与知识检索}]

团队在开赛初期确定本赛题为基于宏基因组数据的疾病预测问题后，先自行梳理了三类疾病（结直肠癌、炎症性肠病、肥胖症）的已知微生物关联与宏基因组建模惯例，随后要求 AI 检索三疾病宏基因组机器学习预测的文献基准，并整理数据来源队列与已知标志物，作为后续建模的领域锚点。

\textbf{提示词：} 检索三疾病（结直肠癌 / 炎症性肠病 / 肥胖症）宏基因组机器学习预测的文献基准与已知微生物标志物，整理各数据来源队列（Zeller 2014、MetaHIT、Chatelier 2013）的建模惯例与预期性能区间，输出领域知识摘要。

\textbf{AI回复（摘要）：} AI 返回三疾病文献摘要：CRC 宏基因组 ML AUC 常见 0.80--0.90（跨队列衰减至 0.70--0.80）、IBD 10-species 外部验证 AUC $\sim$0.85、肥胖预测 AUC 通常 0.65--0.75（弱于 CRC/IBD）；整理出 Fusobacterium nucleatum 等共识标志物与成分数据 CLR、高维小样本、跨队列泛化衰减等建模惯例。团队核验后确认三数据集性能预期排序 CRC $>$ IBD $>$ Obesity，作为 S1 结果解读的领域基准。

\item[\textbf{交互记录 [AI-1-2] —— S1 方案设计}]

团队判定 S1 为预测类（分类）问题后，确定以 AUC 为主指标、F1/Recall(少数类) 为辅，并选定主模型 Logistic(L2)+CLR（成分数据定和偏相关需 CLR 解除、L2 抗 $n\ll p$ 过拟合）、对照随机森林（原始丰度）、基线单特征最佳阈值+Dummy，统一分层 5 折 CV + LOOCV 兜底。团队将这一方案框架交给 AI 形式化并核对评估协议。

\textbf{提示词：} 按团队确定的方案形式化 S1 建模框架：主模型 Logistic(L2)+CLR（近全零过滤 1331$\to$264 后 CLR 变换，class\_weight='balanced'）、对照 Random Forest（原始丰度）、基线单特征最佳阈值+Dummy；评估协议为分层 5 折 CV（seed=42）+ LOOCV 兜底，主指标 AUC。核对各模型超参与评估口径。

\textbf{AI回复（摘要）：} AI 确认方案自洽并给出实现要点：CLR 零值乘法替换 $\delta=6.5\times10^{-6}$、LogisticRegression(penalty='l2', C=1.0, class\_weight='balanced')、RandomForestClassifier(n\_estimators=500)、StratifiedKFold(5, seed=42)。团队审阅后确认方案锁定，进入代码实现。

\item[\textbf{交互记录 [AI-1-3] —— S1 建模实现}]

团队将 S1 方案交给 AI 实现为完整代码：数据加载、近全零过滤、CLR 变换、主模型与对照模型训练、基线计算、分层 5 折 CV 与 LOOCV 评估，并落盘结果。

\textbf{提示词：} 按 S1 方案实现完整建模代码：加载清洗后数据，近全零过滤（零值占比 $>$95\% 剔除，1331$\to$264），CLR 变换（$\delta=6.5\times10^{-6}$），训练 Logistic(L2)+CLR 与 Random Forest（原始丰度），计算单特征最佳阈值与 Dummy 基线，分层 5 折 CV + LOOCV 评估，输出 AUC/ACC/F1/Recall(少数类) 并落盘 S1-results.pkl。

\textbf{AI回复（摘要）：} AI 完成代码实现并返回结果：三数据集 L2(CLR) AUC 为 Zeller 0.791、metahit 0.887、Chatelier 0.650；RF 在 Zeller/metahit 更优（0.845/0.904）、Chatelier 略优（0.660）；LOOCV 与 5 折 CV 差距均 $<$0.025，模型稳定。团队审阅代码与结果后确认实现正确，进入结果分析。

\item[\textbf{交互记录 [AI-1-4] —— S1 结果分析与归因}]

团队基于 S1 结果，要求 AI 协助完成三数据集性能对比、与基线增益计算、跨疾病差异四重归因（样本量/类别平衡/信号强度/批次效应），并对待裁定项（过拟合判定、small\_adenoma 口径、Chatelier 弱信号）给出裁决依据。

\textbf{提示词：} 基于 S1-results.pkl 完成结果分析：(1) 三数据集 L2/RF 性能对比与相对基线增益；(2) 跨疾病差异按样本量、类别平衡、信号强度、批次效应四重来源归因；(3) 裁决过拟合判定（采用 CV vs LOOCV 口径）、small\_adenoma 主口径（维持归健康）、Chatelier 弱信号诚实标注。

\textbf{AI回复（摘要）：} AI 返回分析：metahit 增益最大（L2 +0.072）、Chatelier 增益最小（+0.010，接近领域下界）；性能差异主要由信号强度与类别平衡主导；过拟合判定采用 CV vs LOOCV 口径（差距 $<$0.025 无过拟合），small\_adenoma 维持归健康主口径。团队核验后确认结论，写入报告。

\item[\textbf{交互记录 [AI-2-1] —— S2 方案设计}]

团队判定 S2 为特征选择子任务后，确定采用稳定性选择（Lasso+bootstrap）为主方法，配两路信号（Fisher 存在/缺失 + Wilcoxon 非零丰度，BH-FDR 校正）、共现分析初探、VIP 独立复现佐证。团队将方案交给 AI 形式化并核对统计口径。

\textbf{提示词：} 按团队确定的方案形式化 S2 特征选择框架：近全零过滤（1331$\to$264）$\to$ CLR 前置 $\to$ Lasso+bootstrap 稳定性选择（$\tau=0.5$，$B=100$）$\to$ 两路信号（Fisher/Wilcoxon + BH-FDR m=1331）$\to$ 共现分析 $\to$ VIP 佐证。核对各统计检验口径与多重比较校正规模。

\textbf{AI回复（摘要）：} AI 确认方案并给出实现要点：Lasso 逆正则化 $C=0.1$、bootstrap 频率 $\hat{\pi}_j\ge\tau$ 入选、BH-FDR 按全 1331 特征校正（医疗宁可严格）、VIP 阈值 1.5。团队审阅后确认方案锁定。

\item[\textbf{交互记录 [AI-2-2] —— S2 建模实现}]

团队将 S2 方案交给 AI 实现为完整代码：近全零过滤、CLR、Lasso+bootstrap 稳定性选择、Fisher/Wilcoxon 两路信号、共现分析、VIP 计算，并落盘结果。

\textbf{提示词：} 按 S2 方案实现完整代码：近全零过滤、CLR 变换、Lasso+bootstrap 稳定性选择（$\tau=0.5$，$B=100$，$C=0.1$）、Fisher 精确检验与 Wilcoxon 秩和检验（BH-FDR m=1331）、共现分析（Spearman 相关 + Fisher 共现检验）、PLS-DA VIP，输出每病稳定标志物清单并落盘 S2-results.pkl。

\textbf{AI回复（摘要）：} AI 完成实现并返回：CRC 4 个稳定标志物（3 个 Fisher 显著、全部命中已知标志物）、IBD 4 个（4/4 Fisher 显著）、Obesity 20 个（0 个 FDR 显著，弱信号）；三病 Jaccard 全 0 完全疾病特异；主导信号几乎全为存在/缺失。团队审阅后确认实现正确。

\item[\textbf{交互记录 [AI-2-3] —— S2 结果分析}]

团队基于 S2 结果，要求 AI 协助完成每病稳定标志物解读、两路信号显著性、共现网络、跨疾病对比，并对待裁定项（C 选择、VIP 佐证、RF 佐证层退化）给出裁决依据。

\textbf{提示词：} 基于 S2-results.pkl 完成结果分析：(1) 每病稳定标志物清单与已知标志物对照；(2) 两路信号显著性（Fisher/Wilcoxon + FDR）；(3) 共现网络与跨疾病 Jaccard 对比；(4) 裁决 C 主口径（推荐 0.1）、VIP 独立复现、RF 佐证层退化处置。

\textbf{AI回复（摘要）：} AI 返回分析：CRC/IBD 各 4 个高置信稳定标志物（CRC 3/4 命中已知标志物，方法有效性锚点强验证）；Obesity 弱信号低可信度；判别信号以存在/缺失为主；RF permutation importance 退化不可用，仅以 VIP 作独立复现佐证。团队核验后确认结论，写入报告。

\item[\textbf{交互记录 [AI-3-1] —— S3 方案设计}]

团队判定 S3 为跨疾病泛化评估问题后，确定采用 leave-one-disease-out（LODO）协议做诚实评估，正式构建并对比四种跨疾病建模策略（直接迁移/共享标志物/分类学聚合/部署校正），并设计紧急回退协议（四策略均值均 $<$0.60 时触发）。团队将方案交给 AI 形式化并核对协议。

\textbf{提示词：} 按团队确定的方案形式化 S3 框架：LODO 协议（留一疾病，3 组合）$\to$ 四策略对比（A 直接迁移 / B 共享标志物 / C 分类学聚合 / D 部署校正）$\to$ 紧急回退协议（四策略均值均 $<$0.60 触发，R1--R4 逐级换模型族，达可用线 0.65 或提升 $\ge$0.10 即交付）。核对防泄漏约束（禁测试集重定阈值）。

\textbf{AI回复（摘要）：} AI 确认方案并给出实现要点：策略 A 用 Logistic(L2)+CLR、策略 B 用共享特征子集重训、策略 C 属级/门级聚合、策略 D 在最优策略上叠加 Platt 校准；回退 R1 树模型 / R2 样本合并 / R3 密度比重加权 / R4 对抗式域适应。团队审阅后确认方案锁定。

\item[\textbf{交互记录 [AI-3-2] —— S3 建模实现}]

团队将 S3 方案交给 AI 实现为完整代码：四策略建模、紧急回退 R1--R4、衰减归因、深度迁移分析、C3 阈值漂移量化，并落盘结果。

\textbf{提示词：} 按 S3 方案实现完整代码：LODO 协议下四策略（A 直接迁移 / B 共享标志物 / C 属级门级聚合 / D 部署校正 Platt 校准）建模对比；四策略均值均 $<$0.60 时触发紧急回退 R1--R4（树模型/样本合并/密度比重加权/对抗式域适应）；衰减归因三分法、深度迁移方向一致性、C3 阈值漂移量化；输出各策略 AUC 并落盘 S3-results.pkl。

\textbf{AI回复（摘要）：} AI 完成实现并返回：四策略均值全部低于 0.60 触发回退；回退链 R1--R4 中 R3 加权域适应最优可达 0.6068，仍低于可用线 0.65；衰减归因显示主导来源为疾病特异信号与标签语义漂移；C3 阈值漂移量化显示 Youden 阈值落测试分布 96.0\% 分位、灵敏度仅 0.024。团队审阅后确认实现正确。

\item[\textbf{交互记录 [AI-3-3] —— S3 结果分析}]

团队基于 S3 结果，要求 AI 协助完成四策略对比、紧急回退证据链、衰减归因、深度迁移、C3 阈值漂移分析，并对待裁定项（Platt 符号约定、域内 AUC 口径、R3 密度比方法）给出裁决依据。

\textbf{提示词：} 基于 S3-results.pkl 完成结果分析：(1) 四策略 $\times$ 3 组合 AUC 对比与回退触发证据；(2) 回退链 R1--R4 结果与最优可达判定；(3) 衰减归因三分法、深度迁移方向一致性、C3 阈值漂移诊断；(4) 裁决 Platt 符号约定、域内 AUC 口径、R3 密度比方法。

\textbf{AI回复（摘要）：} AI 返回分析：四策略均值全部 $<$0.60 触发回退，R3 加权域适应最优可达 0.6068 仍低于可用线 0.65，以完整证据链交付；IBD 衰减最大（$-0.271$）、Obesity 主导标签语义漂移；共享物种方向一致性接近随机（51.2\%，$p=0.536$）。团队核验后确认结论，写入报告。

\end{description}

% ============================================================
% 四、采纳和人工修改情况
% ============================================================
\section*{四、采纳和人工修改情况}

\begin{longtable}{|c|p{5.2cm}|c|p{5.2cm}|}
\hline
序号 & AI生成内容概述 & 是否采纳 & 人工修改情况说明 \\
\hline
[AI-1-1] & 三疾病宏基因组文献基准与已知标志物摘要 & 采纳 & 团队核验文献来源与数据队列，确认性能预期排序后作为领域锚点 \\
\hline
[AI-1-2] & S1 建模框架形式化与实现要点 & 采纳 & 团队确定模型选型与评估协议，AI 仅形式化核对，团队审阅后锁定 \\
\hline
[AI-1-3] & S1 建模代码与结果 & 采纳 & 团队审阅代码与结果，确认实现正确；主口径与过拟合判定由团队裁决 \\
\hline
[AI-1-4] & S1 结果分析与归因 & 采纳 & 团队核验数字与归因逻辑，确认结论后写入报告 \\
\hline
[AI-2-1] & S2 特征选择框架形式化 & 采纳 & 团队确定稳定性选择与两路信号方案，AI 仅形式化核对 \\
\hline
[AI-2-2] & S2 建模代码与结果 & 采纳 & 团队审阅代码与结果，确认实现正确；C 主口径与 RF 退化处置由团队裁决 \\
\hline
[AI-2-3] & S2 结果分析 & 采纳 & 团队核验标志物清单与显著性，确认结论后写入报告 \\
\hline
[AI-3-1] & S3 跨疾病框架形式化 & 采纳 & 团队确定 LODO 协议与四策略、回退协议，AI 仅形式化核对 \\
\hline
[AI-3-2] & S3 建模代码与结果 & 采纳 & 团队审阅代码与结果，确认实现正确；Platt 符号与域内口径由团队裁决 \\
\hline
[AI-3-3] & S3 结果分析 & 采纳 & 团队核验回退证据链与归因，确认结论后写入报告 \\
\hline
\end{longtable}

% ============================================================
% 固定条目
% ============================================================
\section*{附加说明}

\subsection*{[AI-0-0] 终稿编译}

\textbf{所用AI工具名称和版本：}

DeepSeek, DeepSeek-V4, 深度求索（DeepSeek）, 2026-08-22

\textbf{具体使用目的和环节：}

终稿 LaTeX 编译与排版。论文正文（含模型描述、公式推导、结果分析）由参赛队独立撰写，AI 仅将已完成的文字、公式和图表转化为 LaTeX 代码并编译为 PDF，全程不参与内容生成。

\textbf{关键交互记录：}

\begin{quote}
提示词：将以下论文正文编译为 LaTeX，不修改任何实质内容。

AI 回复：[编译完成，输出 PDF]
\end{quote}

\textbf{采纳和人工修改情况：}

AI 编译后人工检查排版效果，修正格式问题，内容未做任何 AI 修改。

\vspace{1em}
\hrule
\vspace{1em}

\subsection*{[AI-0-1] 本文件格式编排}

\textbf{所用AI工具名称和版本：}

DeepSeek, DeepSeek-V4, 深度求索（DeepSeek）, 2026-08-22

\textbf{具体使用目的和环节：}

本文档的条目整理、格式编排与 LaTeX 编译。参赛队整理竞赛期间各阶段的 AI 使用记录后，要求 AI 按照章程附录模板的四部分格式进行排版编译。AI 仅做文本润色与格式编排（章程第八条第 2 项允许范围），不涉及内容生成。

\textbf{关键交互记录：}

\begin{quote}
提示词：按照竞赛附录模板的四部分格式（工具名称版本 / 使用目的和环节 / 关键交互记录 / 采纳和人工修改情况）将 AI 使用记录排版编译为 PDF。

AI 回复：[完成全部记录的排版与编译]
\end{quote}

\textbf{采纳和人工修改情况：}

参赛队审核了全部条目的内容完整性和表述准确性，确认后定稿。

\end{document}
